\begin{frame}{Научная проблема, цель и гипотеза}
\textbf{Проблема.} Определить условия, при которых коннектомное предобучение повышает точность моделирования кривой роста ECAP человека.
\SlideGap
\textbf{Цель.} Разработать метод биофизического моделирования ECAP с предобучением на динамике Drosophila и оценить его дополнительный вклад на независимых данных человека.
\SlideGap
\textbf{Объект.} Формирование и регистрация ECAP в измерительном канале системы стимуляции спинного мозга.
\SlideGap
\textbf{Гипотеза.} Предобучение временного модуля на динамике Drosophila может снизить ошибку кривой роста относительно обучения без переноса и контрольных вариантов при сохранении калибровки неопределённости.
\SlideGap
Проверка должна различить преимущество, его отсутствие, ухудшение и неопределённый результат.
\end{frame}
